\scanpage{29}
\mspage{28}

\begin{center}$-39-$\end{center}\hfill n\textsuperscript{o} 272

Corollaire 1.-- \uline{Pour toute famille} $a_i\in A_i$ $(1\leqslant i\leqslant n)$, \uline{il existe} $x\in B$ \uline{tel que} $x\equiv a_i\ (\mathrm{mod.}\ \underline{\mathrm p}_i)$.

Puisque les $\underline{\mathrm p}_i$ sont des idéaux maximaux de $B$, on a $B_{\underline{\mathrm p}_i}/\underline{\mathrm p}_iB_{\underline{\mathrm p}_i}=B/\underline{\mathrm p}_i$, et l'on peut donc supposer que $a_i\in B$ pour tout $i$. Dans ce cas, le corollaire signifie que l'application $B\longrightarrow\prod B/\underline{\mathrm p}_i$ est surjective, ce qui est un résultat bien connu (cf.Ch.I, \S2, cor. à la prop.4).

Corollaire 2.-- \uline{Il existe des éléments} $x_i\in K$ \uline{tels que} $v_i(x_i)=0$ \uline{et} $v_j(x_i)>0$ \uline{pour} $i\ne j$.

Pour chaque indice $i$, on applique le cor.1 à la famille $a_j=0$ si $j\ne i$, $a_i=1$.

Corollaire 3.-- \uline{Tout anneau de valuation contenant} $B$ \uline{contient l'un des} $A_i$.

Soit $A$ un anneau de valuation contenant $B$, et soit $\underline{\mathrm p}=\underline{\mathrm m}(A)\cap B$; il existe un idéal maximal $\underline{\mathrm p}_i$ de $B$ contenant $\underline{\mathrm p}$, d'où $A_i=B_{\underline{\mathrm p}_i}\subset B_{\underline{\mathrm p}}\subset A$.

[En fait, on a $A=B_{\underline{\mathrm p}}$ en vertu de la prop.1, mais on s'en fout.]

2. \uline{Valuations indépendantes}.

Définition 1.-- \uline{Soient} $A$ \uline{et} $A'$ \uline{deux anneaux de valuation d'un même corps} $K$. \uline{On dit que} $A$ \uline{et} $A'$ \uline{sont indépendants s'il n'existe aucun sous-anneau de} $K$, \uline{distinct de} $K$, \uline{les contenant}.

Bien entendu, on dit que deux valuations $v$ et $v'$ de $K$ sont \uline{indépendantes} si leurs anneaux sont indépendants.

\uline{Exemple}. Si $A$ et $A'$ sont \uline{distincts} et sont tous deux des anneaux de valuation de \uline{rang} $1$, ils sont indépendants, en vertu de la prop.7 du \S4.

Proposition 3.-- (\uline{théorème d'approximation}). \uline{Soient} $v_i$ $(1\leqslant i\leqslant n)$ \uline{des valuations deux à deux indépendantes d'un corps} $K$, \uline{et soient} $\Gamma_i$ \uline{leurs groupes des ordres}. \uline{Soient} $a_i\in K$ \uline{et} $\gamma_i\in\Gamma_i$ $(1\leqslant i\leqslant n)$. \uline{Il existe alors} $x\in K$ \uline{tel que} $v_i(x-a_i)\geqslant\gamma_i$ \uline{pour tout} $i$.
